\section{Weights of logarithmic pluriforms}\label{sec:weights}

The denominator needed in the canonical bundle formula will be detected
by a rational pluriform over a curve.  Its weight compares the corrected
canonical order with the multiplicity of the base parameter.  The
main assertion of this section bounds the denominator using only the
relative dimension and the degree of the pluriform.  Coefficient-one
horizontal boundary will permit induction by residue; the remaining
klt case is proved in Proposition~\ref{prop:klt-weight}.

Let $C$ be a smooth projective complex curve, let $c\in C$, and put
$K=\C(C)$.  A rational uniformizer at $c$ is an element $z\in K$ with
$\ord_c z=1$.  Suppose that $F/K$ is a finitely generated regular field
extension of transcendence degree $s$.  A rational relative
$m$-canonical form is a nonzero element of
$(\bigwedge^s\Omega_{F/K})^{\otimes m}$.  We use the canonical-line
identification given by a fixed wedge order to write its associated
absolute form as
\[
                  \Omega=\phi\wedge(dz/z)^{\otimes m}.
\]
Orders of this absolute pluriform are ordinary canonical orders on
models of $F$ over $\C$.

\begin{definition}\label{def:weight}
For a rational relative $m$-canonical form $\phi$, set
\begin{equation}\label{eq:weight-definition}
 w_z(\phi)=\inf_E
 \frac{\ord_E\bigl(\phi\wedge(dz/z)^{\otimes m}\bigr)+m}
      {m\ord_E z}.
\end{equation}
Here $E$ runs through the normalized divisorial valuations of $F/\C$
whose restriction to $K$ is a positive multiple of $\ord_c$.
\end{definition}

The definition depends only on the field and form.  In the log canonical
setting below the infimum is finite and is attained on a log resolution.

\begin{theorem}[Uniform weight denominator]\label{thm:weight}
For integers $s\ge0$ and $m\ge1$ there is an integer $q(s,m)>0$
with the following property.  Let $C$ be a smooth projective complex
curve, $c\in C$, and $z$ a rational uniformizer at $c$.  Let $V$ be a
normal geometrically integral projective $s$-fold over $K=\C(C)$,
and let $(V,B)$ be a geometrically log canonical pair with effective
rational boundary.  If a rational relative $m$-canonical form $\phi$
satisfies
\begin{equation}\label{eq:weight-log-trivialization}
                         \Div_V(\phi)+mB=0,
\end{equation}
then
\[
                            q(s,m)w_z(\phi)\in\Z.
\]
The integer $q(s,m)$ is independent of the boundary coefficients and
of the curve, variety, and form.
\end{theorem}

Although the boundary is not prescribed in the theorem,
\eqref{eq:weight-log-trivialization} already makes $mB$ integral.
What remains unbounded on an arbitrary model is the multiplicity of a
vertical divisor.  The preparation below makes its contribution visible
in a dlt fibre.

\subsection{Changes of fields and forms}

We first record the valuation rules used throughout the proof.  They
extend the boundary-free rules of \cite[Lemma~4.2]{U}; the boundary does
not enter their proof.

\begin{lemma}\label{lem:weight-rules}
Weights of rational relative pluriforms have the following properties.
\begin{enumerate}
\item They do not depend on the rational uniformizer at the fixed point.
For $a\in K^*$ and $t\ge1$,
\[
 w_z(a\phi)=w_z(\phi)+\frac{\ord_c a}{m},
 \qquad w_z(\phi^{\otimes t})=w_z(\phi).
\]
\item If $F'/F$ is finite and $K$ is algebraically closed in $F'$, then
the pullback of $\phi$ has the same weight over $C$.
\item Let $K'/K$ be finite Galois, let $C'\to C$ be the associated map
of smooth projective curves, and let $c'\in C'$ lie above $c$ with
ramification index $l$.  For a rational uniformizer $u$ at $c'$ and the
pullback $\phi'$ to $FK'/K'$, one has
\begin{equation}\label{eq:weight-base-change}
                              w_u(\phi')=l w_z(\phi).
\end{equation}
\end{enumerate}
\end{lemma}

\begin{proof}
If $z'$ is another uniformizer, $(dz'/z')/(dz/z)$ is a base unit at
$c$.  This proves independence of $z$.  For every valuation in
\eqref{eq:weight-definition},
$\ord_E a=(\ord_c a)\ord_E z$; this gives the first formula, and the
tensor-power formula follows by multiplying numerator and denominator.

For a divisorial valuation $E'$ above $E$ in a finite extension of
characteristic-zero fields, with ramification index $a$, the canonical
ramification formula gives
\begin{equation}\label{eq:weight-finite-orders}
                   \ord_{E'}(\Omega)+m
                      =a\bigl(\ord_E(\Omega)+m\bigr).
\end{equation}
Divisorial valuations restrict and prolong to divisorial valuations.
When the constant field is unchanged, $\ord_{E'}z=a\ord_E z$, so
the individual quotients, and hence their infima, agree.

For a curve extension, the pullback of $dz/z$ is a unit times $du/u$
at $c'$, and
\[
                         \ord_{E'}u=\frac{a}{l}\ord_E z.
\]
Equation~\eqref{eq:weight-finite-orders} therefore multiplies each
quotient by $l$.  Regularity of $F/K$ makes $F$ and $K'$ linearly
disjoint over $K$.  The Galois action on their compositum is thus
available to move a prolongation to the chosen branch $c'$.  Every
valuation downstairs has such a prolongation, proving
\eqref{eq:weight-base-change}.
\end{proof}

\subsection{A model on which the weight is exact}

We adapt the curve preparation in \cite[Section~4]{U4} to obtain an
exact divisor equality on a dlt model in any relative dimension.
The elementary discrepancy calculation behind the preparation will
also be used on semistable models.  Write $A(E;W,\Delta)$ for the log
discrepancy, so a component of coefficient one has discrepancy zero.

\begin{lemma}\label{lem:snc-weight}
Let $f:W\to C$ be a projective model of $F/K$ with $W$ smooth.  Put
$F_W=f^*c$ and $T_W=(F_W)_{\mathrm{red}}$.  Let $H_W$ be an effective
horizontal rational boundary whose coefficients are at most one, and
suppose that $T_W+H_W$ has simple normal crossings near $T_W$.
For $\Omega=\phi\wedge(dz/z)^{\otimes m}$, assume that the horizontal
part of $\Div_W(\Omega)/m+H_W$ is effective.  Then
\begin{equation}\label{eq:weight-snc-minimum}
 w_z(\phi)=\min_{D\subset T_W}
                \frac{\ord_D(\Omega)+m}{m\ord_D z},
\end{equation}
where $D$ ranges over the prime components of $T_W$.
\end{lemma}

\begin{proof}
Denote the minimum on the right by $w$.  After shrinking the curve
around $c$, the rational divisor
\begin{equation}\label{eq:weight-error}
 E_W=\frac1m\Div_W(\Omega)+T_W+H_W-wF_W
\end{equation}
is effective.  For every divisorial valuation over $c$, pullback of
this equality gives
\begin{equation}\label{eq:weight-discrepancy}
 \frac{\ord_E(\Omega)}m+1
   =w\ord_E z+A(E;W,T_W+H_W)+\ord_E E_W.
\end{equation}
The last two terms are nonnegative.  Thus every quotient in
\eqref{eq:weight-definition} is at least $w$.  A component realizing
the displayed minimum gives equality.
\end{proof}

\begin{proposition}\label{prop:prepared}
For the data of Theorem~\ref{thm:weight}, with $s>0$, there is a
projective contraction $f:N\to C$ and an effective horizontal rational
boundary $H$ such that:
\begin{enumerate}
\item $N$ is $\Q$-factorial, its generic fibre is geometrically integral
and birational to $V$, and $(N,T+H)$ is dlt, where
$T=(f^*c)_{\mathrm{red}}$;
\item $K_N+T+H\sim_{\Q,C}0$;
\item for the same field-theoretic form $\Omega$ and
$w=w_z(\phi)$, there is an equality of rational divisors near $T$,
\begin{equation}\label{eq:prepared-equality}
                    \frac1m\Div_N(\Omega)+T+H=w f^*c.
\end{equation}
\end{enumerate}
In particular, on the generic fibre,
$\Div(\phi)+mH_\eta=0$.  If $H$ has no coefficient-one component,
then $(N_\eta,H_\eta)$ is geometrically klt and $K_{N_\eta}$ is
$\Q$-Cartier.
\end{proposition}

\begin{proof}
Choose a smooth projective model $W\to C$ whose generic fibre is a
log resolution of $V$.  Resolve also the closures of the horizontal
boundary and exceptional divisors, together with the fibre over $c$.
Let $B^c_{W_\eta}$ be the crepant subboundary on the generic resolution,
and take $H_W$ to be the closure of its positive part.  Its coefficients
are at most one, and the horizontal divisor
\[
       \frac1m\Div_W(\Omega)+H_W
\]
is effective and exceptional over $V$ on the generic fibre.  The
markings can be chosen so that $(W,T_W+H_W)$ is log smooth.
Lemma~\ref{lem:snc-weight} computes $w$ and gives the effective error
$E_W$ in \eqref{eq:weight-error} near $c$.

On the generic fibre, $K_W+T_W+H_W$ is rationally linearly equivalent
to an effective exceptional divisor over the normal projective variety
$V$.  Its Kodaira dimension is zero: a rational function whose poles
are supported on that exceptional divisor descends to a regular
function on $V$, hence is constant.  In particular it has a section
in a sufficiently divisible degree.  Lemma~\ref{lem:curve-mmp}
therefore gives a terminating MMP over $C$ for $K_W+T_W+H_W$.
Write its output as $(N,T+H)$.  No divisors are extracted, so the
transform of $T_W$ is precisely $(f^*c)_{\mathrm{red}}$.
The generic function field is unchanged and is regular over $K$;
Stein factorization consequently makes $f:N\to C$ a contraction.

Discrepancy comparison preserves the divisible generic-fibre section
spaces through these steps.  Thus the generic adjoint on $N$ still
has Kodaira dimension zero.  Its relative semiample contraction has
zero-dimensional image on the generic fibre.  The image is therefore
a normal curve finite over $C$.  Connectedness of the fibres makes
its function field equal to $K$, so this curve is $C$.  This proves
$K_N+T+H\sim_{\Q,C}0$.

Pushforward of \eqref{eq:weight-error} now gives, near the marked fibre,
\begin{equation}\label{eq:prepared-error}
             \frac1m\Div_N(\Omega)+T+H=w f^*c+E_N,
             \qquad E_N\ge0.
\end{equation}
The restriction of $E_N$ to the generic fibre is effective and
$\Q$-linearly trivial, hence zero.  Thus $E_N$ is vertical, and it
is $\Q$-linearly trivial over $C$.  Such a divisor is locally a
rational multiple of a fibre.  Indeed, a relative principal
trivialization has vertical divisor, so its defining rational function
has neither zeros nor poles on the normal projective generic fibre.
It belongs to $K^*$, because that fibre has no new global functions.
Consequently $E_N=a f^*c$ near $c$ for some rational number $a$.

Applying the discrepancy calculation
\eqref{eq:weight-discrepancy} to \eqref{eq:prepared-error} shows that
all valuation quotients are at least $w+a$: the pair $(N,T+H)$ is lc.
Each component of $T$ has coefficient one, so its quotient is exactly
$w+a$.  The valuation definition still gives $w$, since the function
field and form have not changed.  Hence $a=0$, proving
\eqref{eq:prepared-equality}.

The generic restriction of that equality gives the claimed
trivialization.  A dlt pair with no coefficient-one boundary is klt;
generic restriction, checked on a resolution, preserves this statement
after algebraic closure in characteristic zero.  Finally,
$\Q$-factoriality of $N$ makes $K_{N_\eta}$ $\Q$-Cartier.
\end{proof}

\subsection{Residue and the reduction to klt fibres}

Suppose the prepared model has a coefficient-one horizontal component.
Residue lowers the relative dimension, but that component may acquire
new constants.  The arithmetic Stein-degree theorem controls precisely
this finite extension.

\begin{lemma}\label{lem:residue-reduction}
Let $f:(N,T+H)\to C$ be as in Proposition~\ref{prop:prepared}, with
relative dimension $s\ge1$, and suppose that $S$ is a coefficient-one
component of $H$.  Factor $S\to C$ as
\[
                    S\xrightarrow{h}C_S\xrightarrow{\nu}C,
\]
where $h$ is a contraction and $\nu$ is finite.  Then $C_S$ is a
smooth projective curve and
\[
                              \deg\nu\le D_s
\]
for an integer $D_s$ depending only on $s$.  For every $c'\in C_S$
above $c$, with ramification index $e$, there are data satisfying
Theorem~\ref{thm:weight} in relative dimension $s-1$ and degree $m$
over $\C(C_S)$, with a rational form $\phi_S$ whose weight is
\begin{equation}\label{eq:residue-weight}
                             w_u(\phi_S)=e w_z(\phi)
\end{equation}
for a rational uniformizer $u$ at $c'$.
\end{lemma}

\begin{proof}
Dlt adjunction \cite{KM} makes $S$ normal and gives an effective rational
different $\Theta$ such that $(S,\Theta)$ is lc.  Its Stein curve
is normal and hence smooth.  The generic pair $(N_\eta,H_\eta)$ is
normal, projective and lc, has $H^0(\OO_{N_\eta})=K$, and has
$K_{N_\eta}+H_\eta\simq0$.  Its prime boundary component $S_\eta$
has coefficient one.  Theorem~\ref{thm:stein-degree}, applied with
coefficient lower bound $1$, bounds the relative algebraic closure
of $K$ in $K(S_\eta)$.  This field is exactly $\C(C_S)$, so it gives
the asserted $D_s$.

At the generic point of $S$, the absolute form $\Omega$ has a
logarithmic pole of order $m$.  Choose a local equation $x$ of $S$ and
a rational $s$-form $\beta$, regular at the generic point of $S$, whose
restriction generates its canonical line.  Write
\[
 \Omega=a\bigl(dx/x\wedge\beta\bigr)^{\otimes m},
 \qquad a\text{ a unit along }S.
\]
Its degree-$m$ Poincar\'e residue is the rational $m$-canonical form
\[
                 \Omega_S=(a|_S)(\beta|_S)^{\otimes m}.
\]
This construction is independent of the local equation and commutes
with tensor powers.  Adjunction and
\eqref{eq:prepared-equality} give the equality of actual divisors
\begin{equation}\label{eq:residue-divisor}
                 \frac1m\Div_S(\Omega_S)+\Theta
                           =w(f|_S)^*c
\end{equation}
near the marked fibre.  To verify the identity in this exact degree,
take a tensor power for which the ambient adjunction is Cartier and
apply the pluricanonical adjunction isomorphism.  The resulting divisor
identity is that same positive multiple of \eqref{eq:residue-divisor}, and can
be divided by it.  This uses the rational degree-$m$ residue already
defined at the generic point; it does not require a uniform Cartier
index along $S$.

The generic fibre of $h$ is normal, projective and geometrically
integral.  Its pair induced by $\Theta$ is geometrically lc, as one
sees by generic restriction of a resolution in characteristic zero.
Dividing $\Omega_S$ by $(du/u)^{\otimes m}$ in the canonical-line
identification gives a rational relative form $\phi_S$.  The generic
restriction of \eqref{eq:residue-divisor} is exactly
\[
              \Div(\phi_S)+m\Theta_{\eta_S}=0,
              \qquad \eta_S=\Spec\C(C_S).
\]
It remains to compute the weight at $c'$.

Near $h^{-1}(c')$, the right side of
\eqref{eq:residue-divisor} is $ew h^*c'$.  For every valuation over
$c'$ the corrected order is consequently
\[
       \frac{\ord_E(\Omega_S)}m+1
          =ew\ord_Eu+A(E;S,\Theta).
\]
Log canonicity gives $w_u(\phi_S)\ge ew$.  The fibre of $h$ has a
prime divisor, and any such divisor is a component of the intersection
of $S$ with $T$.  At its generic point the two coefficient-one
branches of the dlt pair are simple normal crossings; adjunction gives
coefficient one in $\Theta$.  The quotient for that divisor is $ew$,
which proves \eqref{eq:residue-weight}.
\end{proof}

\begin{proof}[Proof of Theorem~\ref{thm:weight}]
When $s=0$, geometric integrality gives $V=\Spec K$, and
$\phi\in K^*$.  The definition gives
$w_z(\phi)=\ord_c(\phi)/m$, so $q(0,m)=m$ works.

Proceed by induction on $s$.  Use Proposition~\ref{prop:prepared}.
If its horizontal boundary has a coefficient-one component,
Lemma~\ref{lem:residue-reduction} and the induction hypothesis give
\[
                       q(s-1,m)e w_z(\phi)\in\Z
                       \quad\hbox{for some }1\le e\le D_s.
\]
Thus $q(s-1,m)\operatorname{lcm}(1,\ldots,D_s)$ clears the weight
in this case.  Otherwise the prepared generic pair is geometrically
klt with $\Q$-Cartier canonical divisor.  Proposition~\ref{prop:klt-weight},
proved independently of this induction below, gives a clearing integer
depending only on $s,m$.  Taking a common multiple of these two
integers completes the induction.
\end{proof}
